% Hamiltonian with arrows to its model-specific quantities (tikzmark; two passes).
% H, the stage index, diag, the dagger and theta stay unexplained by request.
\noindent\makebox[\linewidth][c]{\fontsize{54}{64}\selectfont
	\(\displaystyle H_\ell(\tikzmarknode{C}{C})=
	\tikzmarknode{Q}{\color{accent}Q_\ell}\,
	\operatorname{diag}\bigl(
	\tikzmarknode{e}{\color{accent}e_\ell}+\tikzmarknode{s}{\color{csred}s_\theta(C)}
	\bigr)\,{\color{accent}Q_\ell^\dagger}\)}%
\begin{tikzpicture}[remember picture,overlay,>=Latex,
		note/.style={font=\sffamily\fontsize{27}{34}\selectfont,inner sep=0pt,align=center},
		pointer/.style={->,line width=2pt,shorten <=3.5mm,shorten >=2.5mm}]
	\node[note,text=accent,anchor=south] (lQ) at ([yshift=20mm]Q.north)
	{learned basis\\\(1024\times128\)};
	\draw[pointer,accent] (lQ.south) -- (Q.north);
	\node[note,text=accent,anchor=south] (le) at ([yshift=20mm]e.north) {learned\\energies};
	\draw[pointer,accent] (le.south) -- (e.north);
	\node[note,text=navy,anchor=north] (lC) at ([yshift=-20mm]C.south)
	{context: up to\\128 past tokens};
	\draw[pointer,navy] (lC.north) -- (C.south);
	\node[note,text=csred,anchor=north] (ls) at ([yshift=-20mm]s.south)
	{context shifts from\\classical pooling};
	\draw[pointer,csred] (ls.north) -- (s.south);
\end{tikzpicture}
